\documentclass{article}
\linespread{1.3}
\usepackage[margin=50pt]{geometry}
\usepackage{amsmath, amsthm, amssymb, amsthm, tikz, fancyhdr, float, graphicx, spverbatim, systeme}
\pagestyle{fancy}
\renewcommand{\headrulewidth}{0pt}
\newcommand{\changefont}{\fontsize{15}{15}\selectfont}

\fancypagestyle{firstpageheader}{
  \fancyhead[R]{
    \changefont
    \parbox[t]{4cm}{ % Adjust width as needed
      Michael Huang\\
      EN.555.644.81\\
      Homework 4
    }
  }
}

\begin{document}

\thispagestyle{firstpageheader}
{\Large 

\section*{5.2.}
What is the difference between the forward price and the value of a forward contract? \\ \\


}
{\Large 

\section*{5.4.}
A stock index currently stands at 350. The risk-free interest rate is 8\% per annum (with continuous compounding) and the dividend yield on the index is 4\% per annum. What should the futures price for a four-month contract be? \\ \\


}
{\Large 

\section*{5.6}
Explain carefully the meaning of the terms convenience yield and cost of carry. What is the relationship between futures price, spot price, convenience yield, and cost of carry? \\ \\


}
{\Large 

\section*{5.7}
Explain why a foreign currency can be treated as an asset providing a known yield. \\ \\


}
{\Large 

\section*{5.12}
Suppose that the risk-free interest rate is 10\% per annum with continuous compounding and that the dividend yield on a stock index is 4\% per annum. The index is standing at 400, and the futures price for a contract deliverable in four months is 405. What arbitrage opportunities does this create? \\ \\


}
{\Large 

\section*{5.16}
Suppose that \(F_1\) and \(F_2\) are two futures contracts on the same commodity with times to maturity, \(t_1\) and \(t_2\), where \(t_2 > t_1\). Prove that 
\[
  F_2 \leq F_1 e^{r(t_2 - t_1)}
\]
where \(r\) is the interest rate (assumed constant) and there are no storage costs. For the purposes of this problem, assume that a futures contract is the same as a forward contract. \\ \\


}
{\Large 

\section*{5.17.}
When a known future cash outflow in a foreign currency is hedged by a company using a forward contract, there is no foreign exchange risk. When it is hedged using futures contracts, the daily settlement process does leave the company exposed to some risk. Explain the nature of this risk. In particular, consider whether the company is better off using a futures contract or a forward contract when... \\ \\
Assume that the forward price equals the futures price

\subsection*{a)} 
The value of the foreign currency falls rapidly during the life of the contract \\ \\


\subsection*{b)}
The value of the foreign currency rises rapidly during the life of the contract \\ \\


\subsection*{c)}
The value of the foreign currency first rises and then falls back to its initial value \\ \\


\subsection*{d)}
The value of the foreign currency first falls and then rises back to its initial value \\ \\

}
{\Large 

\section*{5.20}
Show that equation (5.3) is true by considering an investment in the asset combined with a short position in a futures contract. Assume that all income from the asset is reinvested in the asset. Use an argument similar to that in footnotes 2 and 4 of this chapter and explain in detail what an arbitrageur would do if equation (5.3) did not hold. \\ \\


}
{\Large 

\section*{5.27}
An index is 1,200. The three-month risk-free rate is 3\% per annum and the dividend yield over the next three months is 1.2\% per annum. The six-month risk-free rate is 3.5\% per annum and the dividend yield over the next six months is 1\% per annum. Estimate the futures price of the index for three-month and six-month contracts. All interest rates and dividend yields are continuously compounded. \\ \\ 


}

\end{document}